\documentclass[11pt]{article}
\usepackage[utf8]{inputenc}
\usepackage[T1]{fontenc}
\usepackage{geometry}
\usepackage{enumitem}
\usepackage{xcolor}
\usepackage{amssymb}
\usepackage{booktabs}
\usepackage{array}
\usepackage{longtable}

\geometry{margin=0.75in}

\title{\textbf{Midterm Assignment: Grading Guide}\\
\large Instructor Reference for Objective Assessment}
\author{Prompt Engineering Course}
\date{}

\begin{document}

\maketitle

\section*{Grading Philosophy}

This rubric is designed for \textbf{objective, efficient grading} with minimal subjectivity. Each criterion has clear indicators for full credit, partial credit, or no credit.

\textbf{Time estimate:} 15--20 minutes per student submission.

\section*{Quick Grading Workflow}

\begin{enumerate}
    \item Verify submission format (PDF, all parts present) --- deduct 5 points if incomplete
    \item Grade each part using the rubrics below
    \item Check word counts --- deduct 1 point per section if significantly over/under (>20\%)
    \item Verify screenshots are included and readable
    \item Total points and provide brief feedback (2--3 sentences per part)
\end{enumerate}

\hrule
\vspace{0.5cm}

\section{Part 1: Prompt Comparison Analysis (30 points)}

\subsection*{Detailed Rubric}

\begin{longtable}{|p{5cm}|p{2cm}|p{8cm}|}
\hline
\textbf{Criterion} & \textbf{Points} & \textbf{Indicators} \\
\hline
\endfirsthead
\hline
\textbf{Criterion} & \textbf{Points} & \textbf{Indicators} \\
\hline
\endhead

Task description clear and appropriate & 3 & 
\textbf{Full (3):} Domain stated, task is specific and realistic \\
\textbf{Partial (2):} Domain stated but task somewhat vague \\
\textbf{Minimal (1):} Domain stated but task unclear \\
\textbf{None (0):} Missing or inappropriate \\
\hline

Version A demonstrates vague prompting & 4 & 
\textbf{Full (4):} Clearly lacks specificity, context, constraints \\
\textbf{Good (3):} Somewhat vague but has some detail \\
\textbf{Partial (2):} Not clearly distinguishable from Version B \\
\textbf{Minimal (1):} Minimal difference from Version B \\
\textbf{None (0):} Missing \\
\hline

Version B applies $\geq$3 techniques with labels & 9 & 
\textbf{Full (9):} 3+ techniques correctly applied and labeled (role, delimiters, few-shot, CoT, structured output, etc.) \\
\textbf{Good (7--8):} 3 techniques applied, labels may be incomplete \\
\textbf{Partial (5--6):} 2 techniques, labeling unclear \\
\textbf{Minimal (2--4):} 1--2 techniques, techniques not clearly identified \\
\textbf{None (0--1):} No clear improvement or missing \\
\hline

Screenshots included for both versions & 3 & 
\textbf{Full (3):} Both screenshots present and readable \\
\textbf{Partial (2):} Both present but one illegible \\
\textbf{Minimal (1):} One screenshot missing or illegible \\
\textbf{None (0):} No screenshots \\
\hline

Analysis identifies 3+ concrete improvements & 5 & 
\textbf{Full (5):} 3+ specific, measurable improvements identified (e.g., ``Version B output was 200 words vs.\ 50 words'', ``included specific legal clauses'') \\
\textbf{Good (4):} 3 improvements with some specificity \\
\textbf{Partial (3):} 2 improvements or vague descriptions \\
\textbf{Minimal (1--2):} 1 improvement identified \\
\textbf{None (0):} No analysis \\
\hline

Analysis explains WHY improvements occurred using course concepts & 6 & 
\textbf{Full (6):} Explicitly references course concepts (role prompting, context, constraints, etc.) and explains causal relationship clearly \\
\textbf{Good (5):} References concepts with clear connection \\
\textbf{Partial (3--4):} References concepts but connection unclear \\
\textbf{Minimal (1--2):} Generic explanation without course terminology \\
\textbf{None (0):} No explanation \\
\hline

\textbf{TOTAL} & \textbf{30} & \\
\hline
\end{longtable}

\subsection*{Common Issues to Watch For}

\begin{itemize}
    \item \textbf{Too similar:} Version A and B are nearly identical (deduct from criterion 2)
    \item \textbf{Mislabeled techniques:} Student claims ``chain-of-thought'' but prompt doesn't show reasoning steps (deduct from criterion 3)
    \item \textbf{Superficial analysis:} Says ``B is better'' without specific evidence (deduct from criterion 5)
    \item \textbf{Missing screenshots:} Common oversight (deduct from criterion 4)
\end{itemize}

\hrule
\vspace{0.5cm}

\section{Part 2: Practical Prompt Design (40 points)}

\subsection*{Detailed Rubric}

\begin{longtable}{|p{5cm}|p{2cm}|p{8cm}|}
\hline
\textbf{Criterion} & \textbf{Points} & \textbf{Indicators} \\
\hline
\endfirsthead
\hline
\textbf{Criterion} & \textbf{Points} & \textbf{Indicators} \\
\hline
\endhead

Framework selection justified appropriately & 5 & 
\textbf{Full (5):} Clear, logical explanation of why framework suits task with specific reasoning \\
\textbf{Good (4):} Good explanation with adequate reasoning \\
\textbf{Partial (2--3):} Justification present but weak reasoning \\
\textbf{Minimal (1):} Framework stated, no justification \\
\textbf{None (0):} Missing \\
\hline

All framework components present and clearly labeled & 10 & 
\textbf{CO-STAR (6 components):} Full (10) = all 6; Good (8) = 5; Partial (6) = 4; Minimal (4) = 2--3 \\
\textbf{RACI (4 components):} Full (10) = all 4; Good (8) = 3; Partial (5) = 2 \\
\textbf{IDEA (4 components):} Full (10) = all 4; Good (8) = 3; Partial (5) = 2 \\
\textbf{SMART (5 criteria):} Full (10) = all 5; Good (8) = 4; Partial (6) = 3; Minimal (4) = 2 \\
\textbf{Labels clear:} Deduct 2 points if components not labeled \\
\hline

Prompt is well-structured, clear, and complete & 5 & 
\textbf{Full (5):} Prompt is coherent, specific, actionable, no ambiguity \\
\textbf{Good (4):} Minor clarity issues \\
\textbf{Partial (3):} Some confusion or missing details \\
\textbf{Minimal (1--2):} Poorly structured or incomplete \\
\textbf{None (0):} Missing \\
\hline

AI output included with annotations & 3 & 
\textbf{Full (3):} Output included, student annotated where AI followed framework \\
\textbf{Partial (2):} Output included, no annotations \\
\textbf{Minimal (1):} Partial output \\
\textbf{None (0):} Missing \\
\hline

Identifies 2+ weaknesses and creates effective refinement prompt & 8 & 
\textbf{Full (8):} 2+ specific weaknesses identified, refinement prompt addresses them directly and effectively \\
\textbf{Good (6--7):} 2 weaknesses, refinement prompt addresses them well \\
\textbf{Partial (4--5):} 2 weaknesses, refinement prompt partially addresses \\
\textbf{Minimal (2--3):} 1 weakness, weak refinement \\
\textbf{Poor (1):} Vague issues, generic refinement \\
\textbf{None (0):} Missing \\
\hline

Improved output demonstrates measurable enhancement & 4 & 
\textbf{Full (4):} Clear, significant improvement visible in output (more specific, better structured, addresses identified weaknesses) \\
\textbf{Good (3):} Noticeable improvement \\
\textbf{Partial (2):} Slight improvement \\
\textbf{Minimal (1):} Minimal or unclear improvement \\
\textbf{None (0):} No discernible improvement or missing \\
\hline

Framework effectiveness reflection is thoughtful and specific & 5 & 
\textbf{Full (5):} Evaluates framework strengths/limitations, considers alternatives, demonstrates learning \\
\textbf{Good (4):} Thoughtful but general \\
\textbf{Partial (2--3):} Superficial reflection \\
\textbf{Minimal (1):} Generic statements \\
\textbf{None (0):} Missing \\
\hline

\textbf{TOTAL} & \textbf{40} & \\
\hline
\end{longtable}

\subsection*{Framework Component Checklist}

Use this to quickly verify components are present:

\textbf{CO-STAR:}
\begin{itemize}
    \item[$\square$] Context --- background info
    \item[$\square$] Objective --- clear goal
    \item[$\square$] Style --- writing style specified
    \item[$\square$] Tone --- emotional quality
    \item[$\square$] Audience --- who it's for
    \item[$\square$] Response --- format/structure
\end{itemize}

\textbf{RACI:}
\begin{itemize}
    \item[$\square$] Role --- AI persona
    \item[$\square$] Audience --- target reader
    \item[$\square$] Context --- situation
    \item[$\square$] Instructions --- clear directives
\end{itemize}

\textbf{IDEA:}
\begin{itemize}
    \item[$\square$] Instruct --- main directive
    \item[$\square$] Detail --- specifications
    \item[$\square$] Example --- sample/format
    \item[$\square$] Adjust --- refinement criteria
\end{itemize}

\textbf{SMART:}
\begin{itemize}
    \item[$\square$] Specific --- precise output requirements
    \item[$\square$] Measurable --- quantifiable elements
    \item[$\square$] Actionable --- practical next steps
    \item[$\square$] Relevant --- aligned with goal
    \item[$\square$] Time-bound --- temporal constraints
\end{itemize}

\hrule
\vspace{0.5cm}

\section{Part 3: Critical Reflection (30 points)}

\subsection*{Detailed Rubric}

\begin{longtable}{|p{5cm}|p{2cm}|p{8cm}|}
\hline
\textbf{Criterion} & \textbf{Points} & \textbf{Indicators} \\
\hline
\endfirsthead
\hline
\textbf{Criterion} & \textbf{Points} & \textbf{Indicators} \\
\hline
\endhead

Skill progression demonstrates clear learning trajectory with specific example & 7 & 
\textbf{Full (7):} Articulates evolution from basic to advanced understanding, provides specific ``aha moment'' with context, explains challenge and solution \\
\textbf{Good (5--6):} Shows progression, includes example but lacks depth \\
\textbf{Partial (3--4):} Generic progression description, vague example \\
\textbf{Minimal (1--2):} Superficial reflection \\
\textbf{None (0):} Missing \\
\hline

Practical application is realistic, detailed, and uses appropriate frameworks & 8 & 
\textbf{Full (8):} Realistic scenario, identifies specific frameworks/techniques with justification, anticipates challenges with solutions \\
\textbf{Good (6--7):} Good scenario, frameworks identified but justification weak, some challenges mentioned \\
\textbf{Partial (4--5):} Generic scenario or frameworks not clearly identified, challenges vague \\
\textbf{Minimal (2--3):} Unrealistic scenario, no framework identification \\
\textbf{None (0--1):} Missing or completely off-topic \\
\hline

Limitations \& ethics discussion shows understanding of AI capabilities and responsible use & 8 & 
\textbf{Full (8):} 2 specific limitations identified and explained, relevant ethical consideration with mitigation strategy \\
\textbf{Good (6--7):} 2 limitations, ethical consideration present but solution generic \\
\textbf{Partial (4--5):} 1--2 limitations, ethical discussion superficial \\
\textbf{Minimal (2--3):} Vague limitations, minimal ethics discussion \\
\textbf{None (0--1):} Missing or incorrect \\
\hline

Future learning goals are specific and relevant & 4 & 
\textbf{Full (4):} Specific goals tied to identified skill gaps with clear direction \\
\textbf{Good (3):} Specific goals but connection to learning unclear \\
\textbf{Partial (2):} Generic goals \\
\textbf{Minimal (1):} Vague aspirations \\
\textbf{None (0):} Missing \\
\hline

Writing is clear, well-organized, and meets word count & 3 & 
\textbf{Full (3):} Professional writing, logical flow, within 400--550 words (±10\%) \\
\textbf{Partial (2):} Minor clarity issues or slightly off word count \\
\textbf{Minimal (1):} Disorganized or significantly off word count \\
\textbf{None (0):} Incomprehensible or missing \\
\hline

\textbf{TOTAL} & \textbf{30} & \\
\hline
\end{longtable}

\subsection*{Week 2 Limitations Reference}

Students should reference concepts like:
\begin{itemize}
    \item Hallucinations (making up facts)
    \item Context window limits
    \item Training data cutoff
    \item Cannot reason like humans
    \item Struggles with math/logic
    \item No real-world grounding
    \item Biases in training data
\end{itemize}

\hrule
\vspace{0.5cm}

\section{Overall Grading Notes}

\subsection*{Automatic Deductions}

\begin{itemize}
    \item \textbf{-5 points:} Missing entire part
    \item \textbf{-3 points:} Multiple word count violations (>20\% over/under)
    \item \textbf{-2 points:} Missing AI tool specification
    \item \textbf{-2 points:} Poor formatting (illegible, no section labels)
    \item \textbf{-1 point:} Missing cover page with student info
\end{itemize}

\subsection*{Bonus Opportunities (Optional)}

\begin{itemize}
    \item \textbf{+2 points:} Exceptional creativity or insight in application
    \item \textbf{+2 points:} Goes significantly beyond requirements (e.g., tests 3+ frameworks, creates comprehensive brand guide)
    \item \textbf{Maximum bonus:} 3 points (total cannot exceed 100)
\end{itemize}

\subsection*{Red Flags for Academic Integrity}

Watch for:
\begin{itemize}
    \item Screenshots that don't match prompts described
    \item Suspiciously perfect outputs (AI-generated reflection about AI use)
    \item Identical submissions from multiple students
    \item Prompts copied from online sources without attribution
\end{itemize}

If suspected, request student to walk through their process in office hours.

\hrule
\vspace{0.5cm}

\section{Efficient Grading Checklist}

For each submission:

\begin{enumerate}
    \item \textbf{Format check (2 min):}
    \begin{itemize}
        \item[$\square$] PDF format
        \item[$\square$] All 3 parts present
        \item[$\square$] Cover page with student info
        \item[$\square$] AI tool specified
    \end{itemize}
    
    \item \textbf{Part 1 (5 min):}
    \begin{itemize}
        \item[$\square$] Task description clear (3 pts)
        \item[$\square$] Version A is clearly vague (4 pts)
        \item[$\square$] Version B has 3+ labeled techniques (9 pts)
        \item[$\square$] Screenshots present (3 pts)
        \item[$\square$] Analysis identifies 3+ improvements (5 pts)
        \item[$\square$] Explains WHY with course concepts (6 pts)
    \end{itemize}
    
    \item \textbf{Part 2 (6 min):}
    \begin{itemize}
        \item[$\square$] Framework justified (5 pts)
        \item[$\square$] All components present and labeled (10 pts) --- use checklist
        \item[$\square$] Prompt clear and complete (5 pts)
        \item[$\square$] Output with annotations (3 pts)
        \item[$\square$] Refinement addresses 2 weaknesses (8 pts)
        \item[$\square$] Improvement visible (4 pts)
        \item[$\square$] Reflection thoughtful (5 pts)
    \end{itemize}
    
    \item \textbf{Part 3 (4 min):}
    \begin{itemize}
        \item[$\square$] Skill progression with example (7 pts)
        \item[$\square$] Practical application with frameworks (8 pts)
        \item[$\square$] 2 limitations + ethics (8 pts)
        \item[$\square$] Future goals (4 pts)
        \item[$\square$] Word count 400--550 (3 pts)
    \end{itemize}
    
    \item \textbf{Final review (2 min):}
    \begin{itemize}
        \item[$\square$] Apply any deductions
        \item[$\square$] Consider bonus (optional)
        \item[$\square$] Total score
        \item[$\square$] Write 2--3 sentence feedback per part
    \end{itemize}
\end{enumerate}

\textbf{Total grading time per student: 15--20 minutes}

\hrule
\vspace{0.5cm}

\section{Sample Feedback Templates}

\subsection*{Part 1 Feedback}

\textbf{Strong submission:}\\
``Excellent comparison. Your Version B effectively used role prompting, few-shot examples, and structured output. Analysis clearly identified improvements (specificity, format, tone) and explained WHY using course concepts. Well done!''

\textbf{Needs improvement:}\\
``Version B shows some improvement but techniques need clearer labeling. Analysis would benefit from more specific examples (e.g., word count, structural differences). Reference Week 3 concepts more explicitly.''

\subsection*{Part 2 Feedback}

\textbf{Strong submission:}\\
``CO-STAR framework applied comprehensively with all 6 components clearly labeled. Refinement prompt effectively addressed identified weaknesses (specificity, tone). Reflection demonstrates strong framework understanding.''

\textbf{Needs improvement:}\\
``RACI framework incomplete --- missing Context component. Prompt is clear but refinement doesn't address specific weaknesses. Explain WHY framework was chosen for this task.''

\subsection*{Part 3 Feedback}

\textbf{Strong submission:}\\
``Thoughtful reflection showing clear progression from basic to advanced prompting. Practical application is realistic with appropriate framework choices. Ethics discussion considers hallucination risks and proposes verification steps.''

\textbf{Needs improvement:}\\
``Reflection is general --- provide specific `aha moment' example. Practical application needs concrete framework identification (which Week 3/4 technique?). Discuss LLM limitations from Week 2 more explicitly.''

\end{document}
